\section{Воспроизводство основного капитала в качестве ключевого фактора промышленного цикла}

\subsection{Специфическая роль основного капитала}
По мере развития научно-технического прогресса капиталистическое производство становится все более технологически развитым, рабочий из самостоятельного производителя постепенно превращается в придаток машины, а его труд становится все более механизированным и опосредуется сложной системой машин и оборудования, производственных зданий и сооружений, а также транспортных средств, вследствие чего основной капитал приобретает все более значимую роль в процессе общественного воспроизводства, обуславливая тем самым рост органического строения капитала.

Машины и оборудование, представляя собой основной капитал, составляют особую часть постоянного капитала, не потребляемую полностью в пределах единственного производственного цикла в отличие от сырья, материалов, топлива и иных элементов оборотного капитала, а функционирующую в течение продолжительного периода времени посредством многократного участия в производстве; при этом сохраняющую свое вещественное содержание и подвергающуюся физическому и моральному износу в ходе постепенного переноса собственной стоимости на стоимость готовой продукции, неразрывно связывающего процесс капиталистического воспроизводства с длительным оборотом авансированной стоимости и необходимостью ее периодического полного возмещения в натуральной форме.

Поскольку основной капитал, превращаясь в преобладающий элемент структуры общественного производства лишь на определенной исторической стадии развития капитализма, связанной со становлением крупной машинной индустрии, углублением общественного разделения труда, расширением масштабов кооперации и усилением взаимозависимости отдельных отраслей, создает материально-техническую основу массового производства, постольку кризисы общего перепроизводства и промышленный цикл начинают проявляться не как случайные и внешние нарушения хозяйственной динамики, а как закономерная, периодически воспроизводящаяся форма движения капиталистической экономики.

\subsection{Инертность воспроизводства основного капитала}
Преобладание основного капитала в структуре общественного производства придает капиталистическому воспроизводству более инертный характер, поскольку машины, оборудование, здания, сооружения и иные долговременно функционирующие элементы производственного аппарата требуют значительных первоначальных авансирований стоимости, повышают общий срок оборота капитала, увеличивают издержки их обслуживания, ремонта и эксплуатации, а также затрудняют свободное перемещение капитала между отраслями, вследствие чего возникающие в процессе накопления противоречия и диспропорции не могут быть немедленно устранены рыночным переливом капитала, а накапливаются и разрешаются лишь в течение относительно длительного периода.

Инертность воспроизводства основного капитала проявляется не только на уровне отдельного предприятия, но и в масштабах всего общественного производства, структурно включающего в себя два основных подразделения: \textit{\uppercase\expandafter{\romannumeral 1} подразделение}, производящее средства производства, и \textit{\uppercase\expandafter{\romannumeral 2} подразделение}, производящее средства потребления; причем в рамках первого из подразделений особое значение приобретает сфера производства машин и оборудования в совокупности с иными элементами основного капитала, обуславливающая возможность расширенного воспроизводства как самого первого подразделения, так и второго подразделения, поскольку обновление и расширение производственного аппарата в одних отраслях предполагает соответствующее развитие других поставляющих средства производства отраслей.

Подобная взаимосвязь между отраслями и подразделениями общественного производства обуславливает накопительный характер перепроизводства, поскольку избыточное расширение выпуска в одной отрасли создает повышенный спрос на машины, оборудование, сырье и материалы в смежных отраслях, тем самым побуждая их к дополнительному расширению объема производства, тогда как последующее обнаружение ограниченности платежеспособного спроса и невозможности реализации произведенных товаров превращает частичное перепроизводство в цепную реакцию, охватывающую все более широкий круг отраслей и подготавливающую переход от отдельных диспропорций к кризису общего перепроизводства.

\subsection{Периодичность воспроизводства основного капитала}
Обновление основного капитала осуществляется большинством предприятий в начале фазы подъема промышленного цикла, поскольку именно в данный момент общая конъюнктура капиталистического хозяйства оказывается наиболее благоприятной для новых крупных авансирований капитала: снизившаяся в период кризиса общего перепроизводства норма прибыли постепенно восстанавливается к началу фазы подъема за время оживления, возможности сбыта продукции расширяются под воздействием стабилизации платежеспособного спроса, загрузка производственных мощностей возрастает, а увеличение ожидаемых доходов в контексте перспектив дальнейшего экономического роста выступает для капиталистов в качестве стимула расширения и технического перевооружения производства.

Моменты обновления основного капитала на различных предприятиях имеют тенденцию к синхронизации под давлением конкуренции, поскольку каждый индивидуальный капиталист, с одной стороны, стремится к наиболее раннему внедрению новой техники, позволяющей снизить индивидуальные издержки производства, повысить производительность труда и получить временное преимущество над собственными конкурентами, однако с другой стороны, заинтересован в наиболее продолжительном использовании уже функционирующих машин и оборудования в целях обеспечения максимально полного переноса их стоимости на стоимость готовой продукции во избежание преждевременного обесценения еще не возмещенной части авансированного капитала; в связи с чем большинство предприятий выжидает подходящего момента ближе к концу срока физического износа, пока кто-нибудь из капиталистов не решится первым обновить основной капитал, а после незамедлительно приступает к модернизации своих производственных мощностей в целях сохранения собственной конкурентоспособности и положения на рынке.

Таким образом, воспроизводство основного капитала в масштабе капиталистической экономики приобретает периодический и массовый характер, поскольку необходимость его обновления обусловлена объективным сроком физического и морального износа машин и оборудования, а моменты его полного авансирования и последующего возмещения синхронизируются воздействием рыночной конъюнктуры, колебаний нормы прибыли и конкурентной борьбы, причем возникающие от отраслевого ритма отклонения с течением времени нивелируются под давлением вышеперечисленных факторов и не влекут за собой нарушение периодичности данного процесса. $\dagger$ (Приложение \ref{app:model}, стр. \pageref{app:model})

\subsection{Циклообразующая роль основного капитала}
Циклообразующая роль основного капитала, выступающего материально-вещественной основой крупного машинного производства и концентрирующего в себе значительные массы авансированной стоимости, заключается не в простой необходимости его регулярного возмещения, а в обострении имманентных противоречий капиталистического способа производства в ходе процесса его массового обновления, обуславливающего не только периодическую потребность модернизации производственного аппарата, но и само движение промышленного цикла.

Возникающая под давлением физического износа и роста ожидаемых доходов в начале фазы подъема потребность обновления основного капитала стимулирует рост объема выпуска средств производства, влекущий за собой расширение и технологическую модернизацию производственной базы капиталистической экономики в форме новых машин и оборудования, обуславливая дальнейшее увеличение производства средств потребления, вступающее в конфликт с ограниченностью платежеспособного спроса и пролетарским положением народных масс и тем самым подготавливающее предпосылки дальнейшего кризиса общего перепроизводства.

Таким образом, технологическая модернизация в совокупности с расширением производственных мощностей, способствуя росту органического строения капитала и усилению разрыва между растущими возможностями производства и ограниченными рамками капиталистического потребления, формирует предпосылки дальнейшего кризиса и влечет за собой падение нормы прибыли; однако в то же время падение нормы прибыли в ходе кризиса побуждает капиталистов к расширению и модернизации производства в начале следующей фазы подъема в целях компенсации издержек более низкого уровня доходности, в связи с чем обновление основного капитала в конченом итоге одновременно выступает как причиной, так и следствием кризиса общего перепроизводства, обуславливая циклическую форму движения капиталистической экономики.

\subsection{Модифицирующие факторы промышленного цикла} %\textcolor{red}{Черновик}
\textit{Коммерческий и ссудный кредит}, расширяя масштабы реализации товаров и временно преодолевая границы наличного денежного капитала, способен ускорять фазу оживления и подъема, поскольку позволяет капиталистам продолжать производство, закупать сырье, обновлять основной капитал и расширять сбыт еще до фактического возвращения авансированной стоимости, однако тем самым он одновременно углубляет скрытые диспропорции капиталистического воспроизводства, откладывая момент их обнаружения и усиливая разрушительность кризиса, когда цепь взаимных платежей прерывается, доверие между капиталистами подрывается, а необходимость погашения долгов обнажает действительные пределы платежеспособного спроса.

\textit{Торговля и спекуляции}, опосредуя движение товаров между производством и конечным потреблением, могут на определенное время создавать видимость устойчивого расширения рынка, поскольку рост торговых заказов, накопление товарных запасов и повышение цен стимулируют дальнейшее расширение производства, однако, отрывая динамику спроса от реального потребления и превращая ожидание будущей прибыли в самостоятельный мотив хозяйственной активности, они способствуют накоплению избыточных товарных масс, обостряют несоответствие между производством и реализацией и тем самым ускоряют переход от частичных затруднений сбыта к общему кризису перепроизводства.

\textit{Монополизация экономики}, изменяя форму конкурентной борьбы и концентрируя значительные массы капитала в руках крупных объединений, не устраняет циклического характера капиталистического воспроизводства, а лишь модифицирует его проявления, поскольку монополии, обладая возможностью воздействовать на цены, ограничивать выпуск, перераспределять рынки и временно сдерживать падение прибыли, одновременно усиливают неравномерность развития отраслей, затрудняют свободный перелив капитала, поддерживают существование устаревших производственных мощностей и тем самым придают накоплению противоречий более длительный, скрытый и потенциально более разрушительный характер.

\textit{Финансовый капитал}, возникая на основе сращивания промышленного и банковского капитала и концентрируя контроль над движением денежных ресурсов, инвестиций и кредита, усиливает зависимость промышленного цикла от состояния финансовой сферы, поскольку расширение производства все в большей степени опосредуется банковским кредитом, выпуском ценных бумаг, движением фиктивного капитала и ожиданиями будущей доходности, вследствие чего подъем может сопровождаться чрезмерным раздуванием финансовых обязательств и спекулятивных оценок, а кризис перепроизводства – принимать форму финансового потрясения, выражающегося в падении курсов, банкротствах, кредитном сжатии и обесценении капитала.

\textit{Государственная регулирующая политика}, воздействуя на капиталистическое воспроизводство через бюджетные расходы, налоговую систему, денежно-кредитные меры, государственные заказы, поддержку банков и предприятий, а также через различные формы антикризисного вмешательства, способна смягчать отдельные проявления промышленного цикла, отсрочивать обострение противоречий и поддерживать платежеспособный спрос, однако она не устраняет внутренней основы кризисов, поскольку действует в рамках сохранения частнокапиталистического присвоения и потому чаще перераспределяет, временно локализует или переносит противоречия во времени и пространстве, чем ликвидирует их источник.

\textit{Международный характер} капиталистического производства и обращения, выражающийся в развитии мирового рынка, международного разделения труда, трансграничного движения капитала, внешней торговли и глобальных производственных цепочек, придает промышленному циклу более сложную и неравномерную форму, поскольку кризисные явления, возникающие в одной стране или группе отраслей, через торговые, кредитные, валютные и инвестиционные связи распространяются на другие национальные экономики, одновременно позволяя отдельным капиталам временно смягчать внутренние противоречия за счет внешних рынков, но в конечном счете превращая кризис перепроизводства из локального нарушения в явление мирового капиталистического хозяйства.